\documentclass[french,12pt,a4paper]{book}
\usepackage[utf8]{inputenc}
\usepackage[T1]{fontenc}
%\usepackage[french]{babel}
\usepackage{amsmath}
\usepackage{amsfonts}
\usepackage{fancyhdr}
\usepackage{amssymb}
\usepackage{xcolor} % où xcolor selon l'installation
\definecolor{Prune}{RGB}{99,0,60}
\usepackage{mdframed}
\usepackage{multirow} %% Pour mettre un texte sur plusieurs rangées
\usepackage{multicol} %% Pour mettre un texte sur plusieurs colonnes
\usepackage{scrextend} %Forcer la 4eme  de couverture en page pair
%\usepackage{tikz}
\usepackage{graphicx}
\graphicspath{{misc/}} % figure default path
\usepackage[absolute]{textpos} 
\usepackage{colortbl}
\usepackage{array}
\usepackage{geometry}
\usepackage{hyperref}

\usepackage{lipsum} % à retirer!!!

%% added
\usepackage[english, main=french]{babel} % English language hyphenation
%\usepackage[svgnames]{xcolor} % Enabling colors by their 'svgnames'
\usepackage{microtype} % Better typography
\usepackage[inline]{enumitem} % Required for customising lists
\setlist{noitemsep} % Remove spacing between bullet/numbered list elements
\usepackage{sectsty} % Enables custom section titles
%Options: Sonny, Lenny, Glenn, Conny, Rejne, Bjarne, Bjornstrup
\usepackage[Bjornstrup]{fncychap}
\ChNumVar{\fontsize{76}{80}\usefont{OT1}{pzc}{m}{n}\selectfont}
\usepackage{siunitx} % for SI units
\usepackage{mathtools}
\usepackage[title]{appendix} % for appendices
\usepackage{pgfplots} % for plots using data
\usepackage{xstring}

\usepackage{tikz} % figure de façon générale - T. Chevet (la référence tikz)
	\usetikzlibrary{shapes}
	\usetikzlibrary{arrows.meta}
	\usetikzlibrary{calc}
	\usetikzlibrary{babel}
	\usetikzlibrary{patterns}
	\usetikzlibrary{matrix}
\usepackage{circuitikz}

\usepackage{subfig}

\usepackage{cleveref}
	\crefname{figure}{Figure}{Figures}
	\crefname{section}{Section}{Sections}
	\crefname{appendix}{Annexe}{Annexes}
	\crefname{chapter}{Chapitre}{Chapitres}
	\crefname{equation}{}{}
	\crefname{paragraph}{Paragraphe}{Paragraphes}
	\crefmultiformat{equation}{#2(#1)#3}{ et #2(#1)#3}{, #2(#1)#3}{ et #2(#1)#3}

\setcounter{secnumdepth}{4} % to be able to reference paragraphs 

\usepackage{pgffor} %for loop	

\usepackage{tabularx} %for advanced tables

\usepackage{xifthen} %for conditional test

%%

\newcommand{\jacobian}[3]{ %
	\ifthenelse{\isempty{#1}}
	{%
		\boldsymbol{\mathcal{J}_{#2}^{#3}} %
	}
	{%
		\boldsymbol{\mathcal{J}_{#2}}(#1)^{#3} %
	}	
}

\newcommand{\jacobianHat}[3]{ %
	\ifthenelse{\isempty{#1}}
	{%
		\boldsymbol{\hat{\mathcal{J}}_{#2}^{#3}} %
	}
	{%
		\boldsymbol{\hat{\mathcal{J}}_{#2}}(#1)^{#3} %
	}	
}

\newcommand{\jacobianDot}[3]{ %
	\ifthenelse{\isempty{#1}}
	{%
		\boldsymbol{\dot{\mathcal{J}}_{#2}^{#3}} %
	}
	{%
		\boldsymbol{\dot{\mathcal{J}}_{#2}}(#1)^{#3} %
	}	
}

\newcommand{\vect}[1]{
	\boldsymbol{\MakeLowercase{#1}}
}

\newcommand{\matr}[1]{
	\boldsymbol{\MakeUppercase{#1}}
}

\newcommand{\matrCal}[1]{
	\boldsymbol{\mathcal{\MakeUppercase{#1}}}
}

\DeclareMathOperator{\sign}{sgn}

\newcommand{\R}{\rm I\!R}

%compute get angle for arc
\newcommand{\pgfExtractAngle}[3]{%
	\pgfmathanglebetweenpoints{\pgfpointanchor{#2}{center}}
	{\pgfpointanchor{#3}{center}}
	\global\let#1\pgfmathresult
}

\newcommand\modulo[2]{%
	\@tempcnta=#1
	\divide\@tempcnta by #2
	\multiply\@tempcnta by #2
	\multiply\@tempcnta by -1
	\advance\@tempcnta by #1\relax
	\the\@tempcnta%
}

%%%%%%%%%%%%%%
% Fancy chapter customization
%%%%%%%%%%%%%%

%\colorlet{partbgcolor}{gray!30}% shaded background color for parts
%\colorlet{partnumcolor}{gray}% color for numbers in parts
\colorlet{chapbgcolor}{Prune!90}% shaded background color for chapters
\colorlet{chapnumcolor}{black}% color for numbers in chapters

%\renewcommand*\partformat{%
%	\fontsize{76}{80}\usefont{T1}{pzc}{m}{n}\selectfont%
%	\hfill\textcolor{partnumcolor}{\thepart}}
%
%\makeatletter
%\renewcommand*{\@part}{}
%\def\@part[#1]#2{%
%	\ifnum \c@secnumdepth >-2\relax
%	\refstepcounter{part}%
%	\@maybeautodot\thepart%
%	\addparttocentry{\thepart}{#1}%
%	\else
%	\addparttocentry{}{#1}%
%	\fi
%	\begingroup
%	\setparsizes{\z@}{\z@}{\z@\@plus 1fil}\par@updaterelative
%	\raggedpart
%	\interlinepenalty \@M
%	\normalfont\sectfont\nobreak
%	\setlength\fboxsep{0pt}
%	\colorbox{partbgcolor}{\rule{0pt}{40pt}%
%		\makebox[\linewidth]{%
%			\begin{minipage}{\dimexpr\linewidth+20pt\relax}
%				\ifnum \c@secnumdepth >-2\relax
%				\vskip-25pt
%				\size@partnumber{\partformat}%
%				\fi      %
%				\vskip\baselineskip
%				\hspace*{\dimexpr\myhi+10pt\relax}%
%				\parbox{\dimexpr\linewidth-2\myhi-20pt\relax}{\raggedleft\LARGE#2\strut}%
%				\hspace*{\myhi}\par\medskip%
%			\end{minipage}%
%		}%
%	}%
%	\partmark{#1}\par
%	\endgroup
%	\@endpart
%}

\renewcommand\DOCH{%
	\settowidth{\py}{\CNoV\thechapter}
	\addtolength{\py}{-10pt}
	\fboxsep=0pt%
	\colorbox{chapbgcolor}{\rule{0pt}{40pt}\parbox[b]{\textwidth}{\hfill}}%
	\kern-\py\raise20pt%
	\hbox{\color{chapnumcolor}\CNoV\thechapter}\\%
}

\renewcommand\DOTI[1]{%
	\nointerlineskip\raggedright%
	\fboxsep=\myhi%
	\vskip-1ex%
	\colorbox{chapbgcolor}{\parbox[t]{\mylen}{\CTV\FmTi{#1}}}\par\nobreak%
	\vskip 40pt%
}

\renewcommand\DOTIS[1]{%
	\fboxsep=0pt
	\colorbox{chapbgcolor}{\rule{0pt}{40pt}\parbox[b]{\textwidth}{\hfill}}\\%
	\nointerlineskip\raggedright%
	\fboxsep=\myhi%
	\colorbox{chapbgcolor}{\parbox[t]{\mylen}{\CTV\FmTi{#1}}}\par\nobreak%
	\vskip 40pt%
}
\makeatother

\ChTitleVar{\raggedleft\Large\sffamily\color{white}\usefont{OT1}{phv}{}{n}}

%%%%%%%%%%%%%%%%